% Generated draft: requires mathematical and source review.
\documentclass[11pt]{article}
\usepackage[margin=25mm]{geometry}
\usepackage{fontspec}
\setmainfont{Latin Modern Roman}
\usepackage{amsmath,amssymb,mathtools,unicode-math}
\setmathfont{Latin Modern Math}
\usepackage{enumitem,needspace,xurl,hyperref}
\setlength{\parindent}{0pt}
\setlength{\parskip}{5pt}
\setlength{\emergencystretch}{3em}
\newcommand{\N}{\mathbb{N}}
\newcommand{\Z}{\mathbb{Z}}
\newcommand{\Q}{\mathbb{Q}}
\newcommand{\R}{\mathbb{R}}
\newcommand{\C}{\mathbb{C}}
\newcommand{\F}{\mathbb{F}}
\newcommand{\A}{\mathbb{A}}
\newcommand{\PP}{\mathbb{P}}
\newcommand{\E}{\mathbb{E}}
\newcommand{\eps}{\varepsilon}
\newcommand{\dd}{\,\mathrm{d}}
\begin{document}
\section*{UnsolvedMath ID: 700562}
\textbf{Source problem code:} catalogue entry
\subsection{Definitions and mathematical statement}
For a finite set \(X\), let \(F(X)\) be the free group on \(X\). Given words \(r_1,r_2\in F(X)\), write
\[
G=\langle X\mid r_1,r_2\rangle
 =F(X)/\langle\!\langle r_1,r_2\rangle\!\rangle,
\]
where \(\langle\!\langle r_1,r_2\rangle\!\rangle\) is the smallest normal subgroup of \(F(X)\) containing both relators. The word problem for this fixed presentation is soluble if there is an algorithm which, on every input word in \(X\cup X^{-1}\), halts and correctly decides whether the word represents the identity in \(G\). It is insoluble if no such algorithm exists.

\textbf{Question (Yu. I. Merzlyakov; Kourovka Notebook Problem 9.29).} Does there exist a finite set \(X\) and two words \(r_1,r_2\in F(X)\) such that the group \(\langle X\mid r_1,r_2\rangle\) has insoluble word problem?

This formulation uses the usual finite-presentation interpretation of a two-relator group: a group admitting a presentation with two listed defining relators. The supplied wording does not explicitly specify finiteness of the generating set or require that two be the minimum possible number of defining relators. No minimality condition is imposed here; these interpretive qualifications require review against the cited source.

The supplied record calls this a well-known problem and cites the classical theorem of Magnus that the word problem is soluble in one-relator groups as its motivation. Its registry status is supplied as \emph{open}; that status is not independently verified here.
\subsection{Short English statement}
Does some group presented by finitely many generators and two defining relators have an undecidable word problem, meaning that no algorithm always decides whether an input word represents the identity? This draft states the finite-presentation interpretation and leaves the source's unstated presentation conventions for review.
\subsection{Sources}
\begin{itemize}
\item Local UnsolvedMath registry, permanent ID \(700562\), exact source problem code \texttt{catalogue entry}, title \emph{Kourovka Notebook Problem 9.29}; supplied upstream revision \texttt{d371ee917ac9647fbb06cd6fc529deec6f3dda82}. The supplied record labels the problem open.
\item \emph{The Kourovka Notebook (Unsolved Problems in Group Theory)}, Problem 9.29, proposed by Yu. I. Merzlyakov in the 9th issue (1984). The supplied background identifies the 21st issue edition (2026), \texttt{arXiv:1401.0300v48}. The supplied collection description names E. I. Khukhro and V. D. Mazurov as editors.
\item The supplied statement attributes solubility of the word problem for one-relator groups to a classical theorem of Magnus. It supplies no separate bibliographic reference for that theorem. The notebook citation and mathematical status have not been independently checked.
\end{itemize}
\end{document}
